% This version of CVPR template is provided by Ming-Ming Cheng.
% Please leave an issue if you found a bug:
% https://github.com/MCG-NKU/CVPR_Template.

% \documentclass[review]{cvpr}
\documentclass[final]{cvpr}

\usepackage{times}
\usepackage{epsfig}
\usepackage{graphicx}
\usepackage{amsmath}
\usepackage{amssymb}

% Include other packages here, before hyperref.
\usepackage[pagebackref=true,breaklinks=true,colorlinks,bookmarks=false]{hyperref}

\begin{document}

%%%%%%%%% TITLE
\title{Image Processing and Computer Vision (CSC6051/MDS5112)\\Assignment 1 Report}

\author{Course Staff\\
{\tt\small your.email@link.cuhk.edu.cn}}

\maketitle

%%%%%%%%% ABSTRACT
\begin{abstract}
This report presents detailed implementations and results for Task 1 (image affine transformation) and Task 2 (3D cube projection using a pinhole camera). We describe the mathematical formulation, implementation steps, and provide visual results. Task 3 is outlined briefly as it requires GPU training on the EG1800 dataset.
\end{abstract}

%-------------------------------------------------------------------------
\section{Task 1: Image Affine Transformation}

\paragraph{Implementation Details}
We perform scale $s$, rotation $\theta$, and translation $\mathbf{t}=(t_x,t_y)$ around the image center $\mathbf{c}=(c_x,c_y)$. For any pixel coordinate $\mathbf{p}=(x,y)$,
\begin{equation}
\mathbf{p}' = s\,\mathbf{R}(\theta)\,(\mathbf{p}-\mathbf{c}) + \mathbf{c} + \mathbf{t},
\end{equation}
where
\begin{equation}
\mathbf{R}(\theta) =
\begin{bmatrix}
\cos\theta & -\sin\theta \\
\sin\theta & \cos\theta
\end{bmatrix}.
\end{equation}
To implement this in OpenCV, we construct three non-collinear points near the center, transform them by the formula above, and use \texttt{cv2.getAffineTransform} to compute the $2\\times 3$ affine matrix. The image is then generated by \texttt{cv2.warpAffine} with bilinear interpolation and black padding.

\paragraph{Algorithm (pseudo-code)}
\begin{enumerate}
  \item Read image, compute center $\mathbf{c}$.
  \item Create three reference points around $\mathbf{c}$.
  \item Rotate and scale each point; add translation.
  \item Compute affine matrix from the 3-point pairs.
  \item Warp image to output resolution.
\end{enumerate}

\paragraph{Results \& Analysis}
Figure~\ref{fig:affine} shows the input image (left) and the transformed result (right) for $s=0.5$, $\theta=45^{\circ}$, $t_x=50$, $t_y=-50$. The object is visibly scaled down, rotated counterclockwise, and shifted to the right and upward. The use of the image center as the rotation origin prevents unintended translation artifacts during rotation.
\begin{figure}[htb]
  \centering
  \includegraphics[width=0.95\linewidth]{../results/affine_compare.png}
  \caption{Task 1 affine transformation result.}
  \label{fig:affine}
\end{figure}

%-------------------------------------------------------------------------
\section{Task 2: Project 3D Points to Retina Plane}

\paragraph{Implementation Details}
We implement the pinhole camera model. The intrinsics matrix is
\begin{equation}
\mathbf{K} = \begin{bmatrix} \alpha & 0 & c_x \\ 0 & \beta & c_y \\ 0 & 0 & 1 \end{bmatrix}.
\end{equation}
Given a 3D point $\mathbf{x}_c=(X,Y,Z)^\top$, the projected pixel coordinate is
\begin{equation}
\mathbf{x}_s = \pi(\mathbf{K}\mathbf{x}_c) = \left(\frac{u}{w}, \frac{v}{w}\right),
\end{equation}
where $(u,v,w)^\top = \mathbf{K}\mathbf{x}_c$. We apply this to each cube vertex to project all faces.

The cube is defined by 8 corners of a unit cube, centered at the origin, scaled, rotated (Euler angles), and translated to a specified center. Each of the 6 faces is then projected independently to produce a wireframe visualization.

\paragraph{Results \& Analysis}
Figure~\ref{fig:cube} shows the default projection (center at $(0,0,2)$, rotation $[30,50,0]$, $\\alpha=\\beta=70$, $c_x=c_y=64$). The projected faces form a perspective wireframe with vanishing effects toward the image center. Figure~\ref{fig:cube_var} illustrates a variant with a shifted center and different rotation angles, resulting in asymmetric projection and a different apparent orientation.

Increasing focal length reduces the field of view (zoom-in). Moving the cube farther along $Z$ reduces its projected size. Rotations about $x$ and $y$ axes change which faces are visible and the shape of the projected quadrilaterals.
\begin{figure}[htb]
  \centering
  \includegraphics[width=0.6\linewidth]{../results/task2_cube_default.png}
  \caption{Default cube projection.}
  \label{fig:cube}
\end{figure}

\begin{figure}[htb]
  \centering
  \includegraphics[width=0.6\linewidth]{../results/task2_cube_variant.png}
  \caption{Projected cube under different rotation/position.}
  \label{fig:cube_var}
\end{figure}

%-------------------------------------------------------------------------
\section{Task 3: Portrait Segmentation}

\paragraph{Baseline Training (U-Net)}
We use the provided U-Net with cross-entropy loss and Adam optimizer. The training loop logs training loss, validation loss, and validation pixel accuracy for 20 epochs. These metrics should be plotted against epochs to analyze convergence and generalization.

\paragraph{Training Analysis Protocol}
Plot the three curves over epochs. A typical trend is: training loss decreases smoothly, validation loss decreases then plateaus, and pixel accuracy increases then stabilizes. If validation loss rises while training loss keeps falling, overfitting is likely.

\paragraph{Performance Improvements}
Suggested improvements include:
\begin{itemize}
  \item \textbf{Data augmentation}: random horizontal flips, small rotations, and color jittering to improve robustness.
  \item \textbf{Pretrained backbone}: replace the encoder with a pretrained ResNet (e.g., FCN-ResNet50) and fine-tune.
  \item \textbf{Learning rate scheduling}: use ReduceLROnPlateau or CosineAnnealing to stabilize training.
  \item \textbf{Loss variants}: combine Cross-Entropy with Dice loss to handle class imbalance.
\end{itemize}

\paragraph{Face-aware Preprocessing}
A face detector (e.g., OpenCV Haar cascade or MediaPipe Face Detection) is used to locate the face. We crop an expanded bounding box around the face and resize to the input size (e.g., 224$\times$224). This improves performance on images where the portrait is not centered.

\paragraph{SAM2 Zero-shot Comparison}
Use SAM2 to segment portraits with point or box prompts derived from face detection. Visualize at least 5 examples showing input, prompts, and predicted masks, then compare with the trained model qualitatively (boundary sharpness, background leakage, etc.).

\begin{figure}[htb]
  \centering
  \includegraphics[width=0.7\linewidth]{../results/task3_curves_placeholder.png}
  \caption{Placeholder for training curves (replace with actual plots after training).}
\end{figure}

\section{Conclusion}
We provided detailed implementations and visual results for affine transformation (Task 1) and 3D-to-2D projection (Task 2). The key behaviors of scale/rotation/translation and perspective projection are verified by the generated figures. Task 3 requires GPU training on EG1800; once training is executed, the placeholder plots should be replaced with actual curves and qualitative segmentation examples.

{\small
\bibliographystyle{ieee_fullname}
\bibliography{egbib}
}

\end{document}
